% Generated draft: requires mathematical and source review.
\documentclass[11pt]{article}
\usepackage[margin=25mm]{geometry}
\usepackage{fontspec}
\setmainfont{Latin Modern Roman}
\usepackage{amsmath,amssymb,mathtools,unicode-math}
\setmathfont{Latin Modern Math}
\usepackage{enumitem,needspace,xurl,hyperref}
\setlength{\parindent}{0pt}
\setlength{\parskip}{5pt}
\setlength{\emergencystretch}{3em}
\newcommand{\N}{\mathbb{N}}
\newcommand{\Z}{\mathbb{Z}}
\newcommand{\Q}{\mathbb{Q}}
\newcommand{\R}{\mathbb{R}}
\newcommand{\C}{\mathbb{C}}
\newcommand{\F}{\mathbb{F}}
\newcommand{\A}{\mathbb{A}}
\newcommand{\PP}{\mathbb{P}}
\newcommand{\E}{\mathbb{E}}
\newcommand{\eps}{\varepsilon}
\newcommand{\dd}{\,\mathrm{d}}
\begin{document}
\section*{UnsolvedMath ID: 701516}
\textbf{Source problem code:} catalogue entry
\subsection{Definitions and mathematical statement}
For the finite-system formulation, let \(G\) be a group, let \(m,n\in\Z_{\geq 0}\), and let \(F(x_1,\ldots,x_n)\) be the free group on the unknowns. A system of equations with coefficients in \(G\) is
\[
 w_i(x_1,\ldots,x_n)=1\qquad(1\leq i\leq m),
 \qquad w_i\in G*F(x_1,\ldots,x_n).
\]
For each \(i,j\), let \(a_{ij}\in\Z\) be the total exponent sum of \(x_j\) in \(w_i\); equivalently, it is the image of \(w_i\) under the homomorphism \(G*F(x_1,\ldots,x_n)\to\Z\) that kills \(G\), sends \(x_j\) to \(1\), and sends every other unknown to \(0\). The system is \emph{independent} if its exponent-sum matrix \(A=(a_{ij})\) satisfies
\[
 \operatorname{rank}_{\Q}(A)=m.
\]
An overgroup solution consists of a group \(\bar G\), an injective homomorphism \(\iota:G\hookrightarrow\bar G\), and elements \(h_1,\ldots,h_n\in\bar G\) such that every \(w_i\), evaluated using \(\iota\) on its coefficients and \(x_j=h_j\), equals \(1\).

For a group \(H\), put \(\gamma_1(H)=H\) and \(\gamma_{r+1}(H)=[\gamma_r(H),H]\), and put \(H^{(0)}=H\) and \(H^{(r+1)}=[H^{(r)},H^{(r)}]\). Here \([K,L]\) denotes the subgroup generated by the commutators of elements of \(K\) and \(L\). The group \(H\) is nilpotent if \(\gamma_{c+1}(H)=1\) for some finite \(c\geq0\), and solvable if \(H^{(d)}=1\) for some finite \(d\geq0\).

The three assertions requested in the supplied record are:
\begin{enumerate}
\item[(a)] \textbf{Kervaire--Laudenbach Conjecture (KLC).} For every group \(G\) and every independent system with coefficients in \(G\), there is an overgroup solution.
\item[(b)] \textbf{Nilpotent version.} For every nilpotent group \(G\) and every independent system with coefficients in \(G\), there is an overgroup solution with \(\bar G\) nilpotent.
\item[(c)] \textbf{Solvable version.} For every solvable group \(G\) and every independent system with coefficients in \(G\), there is an overgroup solution with \(\bar G\) solvable.
\end{enumerate}
In (b) and (c), the source specifies no bound on the nilpotency class or derived length of \(\bar G\), and does not require it to equal that of \(G\).

The source additionally attributes to M. Gerstenhaber and O. S. Rothaus the assertion that (a) holds for every locally residually finite group \(G\). Here residually finite means that each nonidentity element has nonidentity image under some homomorphism to a finite group; locally residually finite means that every finitely generated subgroup is residually finite. This result and its attribution are reported from the supplied record, without independent verification.

\emph{Scope qualification.} The supplied statement does not explicitly say whether systems must be finite. The formulation above states the finite-system reading, with the usual rational rank of an integer matrix. Whether the source intends any additional assertion for infinite systems, and what rank convention would apply there, remains unresolved.
\subsection{Short English statement}
An independent finite system of group equations has linearly independent exponent-sum rows over \(\Q\). The three questions ask whether every such system can be solved in an overgroup, with the overgroup required to be nilpotent or solvable when the coefficient group has that property. The source leaves the scope of infinite systems unspecified and labels the record open; that label is supplied metadata, not an independently verified assessment.
\subsection{Sources}
\begin{itemize}
\item Local UnsolvedMath registry: supplied complete record with permanent ID \(701516\), upstream ID \(701516\), and exact source problem code \texttt{catalogue entry}; title: \emph{Kourovka Notebook Problem 18.87}. Supplied upstream revision: \texttt{d371ee917ac9647fbb06cd6fc529deec6f3dda82}. The record was created and updated on 2026-10-16 and carries the status \texttt{open}; this is a supplied registry classification.
\item \emph{The Kourovka Notebook (Unsolved Problems in Group Theory)}, original problem 18.87, 18th issue (2014); the supplied record also identifies the 21st issue edition (2026), \texttt{arXiv:1401.0300v48}. Proposed by V. A. Roman'kov in 2014. The supplied collection description names E. I. Khukhro and V. D. Mazurov as editors of the collection covering Issues 1--21 (1965--2026). These bibliographic details are supplied attribution and have not been independently checked.
\item The supplied statement credits M. Gerstenhaber and O. S. Rothaus with the locally residually finite case of (a). No separate publication citation, theorem locator, or reference link for this result is supplied.
\end{itemize}
\end{document}
